% FILE: 05_evidence_receipts_and_gates.tex
% Project: adversarial
% Thesis Role: adversarial-testing-and-hostile-assumption-verification

\section{Evidence, Receipts, and Gates}
\label{sec:adversarial-evidence}

\subsection{Tests}
\label{sec:adversarial-evidence-tests}

The adversarial corpus contains ten executable scenarios with Python \texttt{assert}
statements in \texttt{adversarial/simulate\_laundering\_spoofing.py}. Each scenario tests a
specific attack vector against the \texttt{IngestionBoundary} and asserts the expected
verdict and failed rule.

\textbf{Scenario 0 (Conforming Ingestion):} Two conforming traces manufactured under the
genuine system key are admitted. The \texttt{assert} statement verifies \texttt{verdict ==
ACCEPTED}.

\textbf{Scenario 1 (Timestamp Laundering):} A conforming trace has the
\texttt{Invoice\_Customer} event timestamp shifted backward by one hour without updating
signatures. The ingestion boundary catches the payload modification via
\texttt{check\_signature\_authenticity} before the monotonicity check fires, because the
payload has changed. The \texttt{assert} verifies \texttt{verdict == REJECTED} and
\texttt{rule\_failed == check\_signature\_authenticity}.

\textbf{Scenario 2 (Event Deletion):} The \texttt{Ship\_Goods} event (index 2) is removed
from a conforming trace. The missing event breaks the hash chain at index 2, because the
hash link of what was formerly index 3 no longer validates against the current predecessor.
The \texttt{assert} verifies \texttt{rule\_failed == check\_cryptographic\_chain}.

\textbf{Scenario 3 (Signature Forgery):} A complete trace is manufactured under
\texttt{FORGED\_SYSTEM\_KEY}, with internally consistent signatures and hash chain links.
The ingestion boundary uses \texttt{hmac.compare\_digest} against the genuine key and
detects the forged signatures. The \texttt{assert} verifies
\texttt{rule\_failed == check\_signature\_authenticity}.

\textbf{Scenario 4 (Raw DataFrame Injection):} A raw \texttt{pd.DataFrame} is submitted as
the log input. The boundary immediately rejects it with \texttt{rule\_failed ==
reject\_unhashed\_dataframe} without computing any hash or inspecting any signature.

\textbf{Scenario 5 (Impossible Velocity):} A trace is manufactured with 100 nanoseconds
between events 0 and 1, with all downstream signatures and chain links recalculated under the
genuine key (so the boundary must evaluate the velocity check independently of the signature
check). The \texttt{assert} verifies \texttt{rule\_failed == check\_impossible\_velocity}.

\textbf{Scenarios 6 through 9} require \texttt{libwasm4pm.dylib} and test FFI boundary
forgery (Ed25519 forged signature returns \texttt{ERR\_CONFORMANCE\_VIOLATION = 0xFB03}),
unstructured data rejection (\texttt{wasm\_parse\_and\_query} on a plain text buffer returns
\texttt{ERR\_CONFORMANCE\_VIOLATION}), out-of-bounds pointer safety (offset
\texttt{0xFFFFFFFF} returns \texttt{ERR\_LIFECYCLE\_VIOLATION = 0xFB05}), and heap
zeroization (sentinel bytes are fully zeroed after \texttt{wasm\_shred\_heap}).

The standards audit script (\texttt{validate\_standards.py}) constitutes an additional
executable test: it exits 0 only when all six placement files satisfy the four structural
requirements.

\subsection{Receipts}
\label{sec:adversarial-evidence-receipts}

The ADVERSARIAL\_RECEIPT in \texttt{receipts/RECEIPT\_REGISTRY.md} is the canonical receipt
for the adversarial corpus. Its fields are:
\begin{itemize}
  \item \texttt{produced\_by}: \texttt{adversarial/} challenge workflow
  \item \texttt{witness}: Chicago TDD hostile assumptions doctrine
  \item \texttt{result\_type}: adversarial challenges to each major claim, with refutation or
        open status
  \item \texttt{criteria\_met}: adversarial artifact count $\geq 3$ (actual: 3); challenges
        cover raw laundering, metric fabrication, and conformance theater
  \item \texttt{gaps\_remaining}: additional adversarial challenges may be warranted for
        M\&A-specific claims
\end{itemize}

The adversarial audit report (\texttt{audits/adversarial\_audit\_v30.1.1.md}) carries a
VERIFIED/PASS verdict with Blake3 checksum over 567 commits across 15 directories and an
Ed25519 signature from \texttt{teamwork\_preview\_worker\_m5\_1}. The audit certifies that
all twelve corpus directories meet or exceed ALIVE\_001 minimum document count thresholds
and that the git commit log maintains a sound, continuous receipt chain.

\subsection{Checkpoints}
\label{sec:adversarial-evidence-checkpoints}

The ALIVE OMEGA checkpoint is declared in
\texttt{checkpoints/PROCESS\_INTELLIGENCE\_ADVERSARIAL\_V30.1.1\_OMEGA.md}. The checkpoint
document states that a twenty-agent adversarial swarm relentlessly tested the boundaries of
the framework; that all attack vectors, recursive logic bombs, and ontological paradoxes were
absorbed and neutralized; and that the framework demonstrates self-healing, adaptive
resilience, and emergent capacity for systemic evolution. The checkpoint is permanent under
the Immutability Doctrine: it cannot be reverted, only extended by a follow-up checkpoint.

The ALIVE\_001 checkpoint (\texttt{checkpoints/PROCESS\_INTELLIGENCE\_ALIVE\_001.md}) is the
baseline gate against which the adversarial corpus was measured. The adversarial audit report
certifies that ALIVE\_001 gate criteria were satisfied at the time of the v30.1.1 audit.

\subsection{ALIVE/PARTIAL/BLOCKED Status}
\label{sec:adversarial-evidence-status}

Status: \textbf{ALIVE}.

The adversarial corpus satisfies all three conditions required for an ALIVE verdict: (1) a
receipt exists (ADVERSARIAL\_RECEIPT in \texttt{receipts/RECEIPT\_REGISTRY.md}); (2) tests
exist (ten executable scenarios in \texttt{simulate\_laundering\_spoofing.py} with
\texttt{assert} statements, plus \texttt{validate\_standards.py}); and (3) a checkpoint
exists (\texttt{checkpoints/PROCESS\_INTELLIGENCE\_ADVERSARIAL\_V30.1.1\_OMEGA.md}, ALIVE
OMEGA). These three conditions are necessary and sufficient for an ALIVE declaration under
the research program's gate criteria. The corpus exceeds the minimum artifact threshold
(actual: 3 core files plus extensive supporting assets).

Open findings (001, 004, 005, 006, 008, 010) do not downgrade the status from ALIVE because
they are research program self-corrections: they identify gaps, notation fixes, and deferred
audits that are actionable and bounded. Finding 001 (graduation bridge one-sided) is
explicitly UNRESOLVABLE PRE-GAP\_001-CLOSE, not a defect in the adversarial corpus itself.
Finding 005 (PM4Py algorithm counts not source-verified) is an audit gap, not an adversarial
corpus defect.
